% Shared preamble for the COMMA 2026 paper.
% This file collects the IOS-compatible parts of the migrated old draft
% preamble and the paper-specific helper macros currently in use.

% Useful packages carried over from the old draft, keeping only items that are
% compatible with the IOS Press class and pdflatex-based compilation.
\usepackage[show]{ed}
\usepackage{latexsym}
\usepackage{amssymb}
\usepackage{amsmath}
\usepackage{amsthm}
\usepackage{booktabs}
\usepackage{graphicx}
\usepackage{xparse}
\usepackage{tikz}
\usepackage{ebproof}
\usepackage{stmaryrd}

% A lightweight non-floating grammar environment. The external
% formal-grammar package clashes with the IOS template's float/caption setup.
\newenvironment{grammar}[1][]{%
  \[
  \begin{array}{@{}l@{\;}c@{\;}l@{\quad}l@{}}
}{%
  \end{array}
  \]
}
\newcommand{\firstcasesubtil}[3]{#1 & ::= & #2 & \textit{#3}\\}
\newcommand{\gralt}{\mid}
\newcommand{\gramtext}[1]{\text{#1}}
\usepackage{cancel}
\usepackage{tabularx}
\usepackage{algorithm}
\usepackage[noend]{algpseudocode}
\usepackage[outline]{contour}
\usepackage{pifont}

\makeatletter
\newcommand\Pimathsymbol[3][\mathord]{#1{\@Pimathsymbol{#2}{#3}}}
\def\@Pimathsymbol#1#2{\mathchoice
  {\@Pim@thsymbol{#1}{#2}\tf@size}
  {\@Pim@thsymbol{#1}{#2}\tf@size}
  {\@Pim@thsymbol{#1}{#2}\sf@size}
  {\@Pim@thsymbol{#1}{#2}\ssf@size}}
\def\@Pim@thsymbol#1#2#3{\mbox{\fontsize{#3}{#3}\Pisymbol{#1}{#2}}}
\input{utxmia.fd}
\DeclareFontShape{U}{txmia}{m}{n}{<->ssub * txmia/m/it}{}
\DeclareFontShape{U}{txmia}{bx}{n}{<->ssub * txmia/bx/it}{}
\newcommand{\pilambdaup}{\Pimathsymbol[\mathord]{txmia}{21}}
\newcommand{\Alambda}{\ensuremath{\tikz[baseline=-.6ex,line width=.1ex]{
    \node[inner sep=0pt] (L) at (0,0.03) {$\pilambdaup$};
    \draw (L.south west)++(.12em,.425ex) -- ++(.2em,0);}}}
\makeatother
\newtheorem{theorem}{Theorem}
\newtheorem{lemma}[theorem]{Lemma}
\newtheorem{corollary}[theorem]{Corollary}
\newtheorem{proposition}[theorem]{Proposition}
\newtheorem{fact}[theorem]{Fact}
\theoremstyle{definition}
\newtheorem{definition}{Definition}
\newtheorem{example}[theorem]{Example}

\newcommand{\matilda}{\ensuremath{\bar{\lambda}\mu\tilde{\mu}}}
\newcommand{\alignarrow}{\mathrel{\mkern30mu\rightarrow\mkern30mu}}
\newcommand{\leftattack}{%
  \tikz[baseline=-0.5ex,scale=0.2]{\draw [thick, arrows={-to[scale=0.5]}] (1.8,0.75) -- (0.7,0.75) -- (1.3,0.25) -- (0,0.25);}%
}
\newcommand{\rightattack}{%
  \tikz[baseline=-0.5ex,scale=0.2]{\draw [thick, arrows={-to[scale=0.5]}] (0,0.25) -- (1.3,0.25) -- (0.7,0.75) -- (1.8,0.75);}%
}
\newcommand{\pro}[1]{\textcolor{black}{\contour{white}{$#1$}}}
\newcommand{\refu}[1]{\textcolor{white}{\contour{black}{$#1$}}}
\newcommand{\contr}[1]{\textcolor{black}{\contour{black}{$#1$}}}
\newcommand{\taut}[1]{\textcolor{white}{\contour{white}{$#1$}}}
\newcommand{\cons}[1]{\textcolor{gray}{\contour{white}{$#1$}}}
\newcommand{\conj}[1]{\textcolor{gray}{\contour{white}{$#1$}}}
\newcommand{\ind}[1]{\textcolor{gray}{\contour{gray}{$#1$}}}

\newcommand{\symcut}[3]{ \hypo{\pro{\Sigma}\vdash #1:\pro{#3}\dashv \refu{\Delta}}
      \hypo{\pro{\Sigma}\vdash \refu{#3}:#2 \dashv \refu{\Delta}}
      \infer2{\pro{\Sigma}\vdash \rangle #1:\contr{#3}:#2\langle \dashv \refu{\Delta}}}
\newcommand{\symaxiom}[2]{%
  \infer0{\pro{\Sigma},#1:\pro{#2}\vdash #1:\pro{#2}\dashv \refu{\Delta}}}
\newcommand{\metasymaxiom}[2]{%
  \infer0{V;\pro{\Sigma},\conj{#1}:\pro{#2}\vdash \conj{#1}:\pro{#2}\dashv \refu{\Delta};E}}
\newcommand{\symproimplintro}[4]{%
  \hypo{\pro{\Sigma},#1:\pro{#2}\vdash #3:\pro{#4}\dashv \refu{\Delta}}
  \infer1{\pro{\Sigma}\vdash \lambda #1:#2.#3:\pro{#2\rightarrow #4}\dashv \refu{\Delta}}}
\newcommand{\symprodiffintro}[4]{%
  \hypo{\pro{\Sigma}\vdash #1:\refu{#2}\dashv \refu{\Delta}}
  \hypo{\pro{\Sigma}\vdash #3:\pro{#4}\dashv \refu{\Delta}}
  \infer2{\pro{\Sigma}\vdash #1\circ #3:\pro{#4-#2}\dashv \refu{\Delta}}}
\newcommand{\symmoixa}[2]{%
  \infer0{\Sigma\vdash \refu{#2}:#1 \dashv \refu{#2}:#1,\refu{\Delta}}}
\newcommand{\metasymmoixa}[2]{%
  \infer0{V;\Sigma\vdash \refu{#2}:\cons{#1} \dashv \refu{#2}:\cons{#1},\refu{\Delta};E}}
\newcommand{\symrefudiffintro}[4]{%
  \hypo{\pro{\Sigma}\vdash \refu{#4}:#3 \dashv \refu{#2}:#1,\refu{\Delta}}
  \infer1{\pro{\Sigma}\vdash \refu{#4-#2}:#3.#1\lambda \dashv \refu{\Delta}}}
\newcommand{\symrefuimplintro}[4]{%
  \hypo{\pro{\Sigma}\vdash #1:\pro{#2}\dashv \refu{\Delta}}
  \hypo{\pro{\Sigma}\vdash \refu{#4}:#3 \dashv \refu{\Delta}}
  \infer2{\pro{\Sigma}\vdash \refu{\sigma\rightarrow\delta}:#1\circ #3:\dashv \refu{\Delta}}}
\newcommand{\symrefuintro}[4]{%
  \hypo{\pro{\Sigma},#1:\pro{#2}\vdash #3:\contr{#4}\dashv \refu{\Delta}}
  \infer1{\pro{\Sigma}\vdash \refu{#2}:#3.#1\mu\dashv \refu{\Delta}}}
\newcommand{\symprointro}[4]{%
  \hypo{\pro{\Sigma}\vdash #3:\contr{#4}\dashv \refu{\Delta},\refu{#2}:#1}
  \infer1{\pro{\Sigma}\vdash \mu #1.#3:\pro{#2}\dashv \refu{\Delta}}}

\newcommand{\envsubpremone}{%
  \hypo{V;\pro{\Sigma}\vdash t[\refu{\sigma}:\alpha] \dashv \refu{\Delta};E}%
}
\newcommand{\envsubpremtwo}{%
  \hypo{V';\pro{\Sigma}\vdash \refu{\sigma}:E \dashv \refu{\Delta};E'}%
}
\NewDocumentCommand{\envsub}{O{\envsubpremone} O{\envsubpremtwo} O{V''} O{\pro{\Sigma}} O{t} O{E} O{\alpha} O{\refu{\Delta}} O{E''}}{%
  \begin{prooftree}
      #1
      #2
      \infer2{#3;#4\vdash #5[#6/#7] \dashv #8;#9}
    \end{prooftree}}
\newcommand{\termsubpremone}{%
  \hypo{V;\pro{\Sigma}\vdash [X:\pro{\sigma}]t \dashv \refu{\Delta};E}%
}
\newcommand{\termsubpremtwo}{%
  \hypo{V';\pro{\Sigma}'\vdash v:\pro{\sigma} \dashv \refu{\Delta}';E'}%
}
\NewDocumentCommand{\termsub}{O{\termsubpremone} O{\termsubpremtwo} O{V''} O{\pro{\Sigma}} O{E} O{\alpha} O{t} O{\refu{\Delta}} O{E''}}{%
  \begin{prooftree}
      #1
      #2
      \infer2{#3;#4\vdash [#5/#6]#7 \dashv #8;#9}
    \end{prooftree}}

\newcommand{\itemR}[2]{%
  \item\leavevmode
  \begin{tabularx}{\linewidth}{@{}Xr@{}}#1 & #2 \\\end{tabularx}%
}
\algblockdefx[Prover Claim]{pclaim}{pdone}[1]{\textbf{P:} I claim #1.}{qed}
\algblockdefx[Skeptic Challenge]{schall}{sdone}[2][$\top$]{\textbf{S:} I challenge you to prove #2 given #1.}{done}
\algblockdefx[Prover start Lemma]{pslem}{psmel}[1]{\textbf{P:} I will first prove a lemma: #1.}{done}
\algblockdefx[Skeptic refute Lemma]{srlem}{srmel}[4][$\top$]{\textbf{S:} Even if so, you still need to prove #2 to meet my challenge #3. I challenge you to prove #4 given #1.}{done}

% AI-authored contribution blocks. Use these to wrap agent-written text so it
% remains visually distinct from MR notes and older draft material.
\makeatletter
\newcommand{\aibeglabel}{BegAIP}
\newcommand{\aiendlabel}{EndAIP}
\newcommand{\aibegmargin}{BAIP}
\newcommand{\aiendmargin}{EAIP}
\newcommand{\aipartlabels}[2]{\def\aibeglabel{#1}\def\aiendlabel{#2}}
\newcommand{\aipartmargins}[2]{\def\aibegmargin{#1}\def\aiendmargin{#2}}
\newenvironment{aipart}[1]%
{\ed@part{}{AI contribution}\aibeglabel\aibegmargin\ifshowednotes\footnotemark\footnotetext{#1}\color[rgb]{0.80,0.25,0.25}\fi}
{\ended@part\aiendmargin}
\makeatother